\begin{table}[htbp]
\centering
\caption{Detection Performance: FL Screen vs.\ Imhof-Style Bid-Level Screens}
\label{tab:imhof_comparison}
\begin{threeparttable}
\small
\begin{tabular}{lccc}
\toprule
Screen & AUC & Data required & DeLong $p$ \\
\midrule
FL screen (this paper) & 0.940 & Participation only & --- \\
Imhof CV only & 0.758 & Bid values & \\
Imhof spread only & 0.789 & Bid values & \\
Imhof composite & 0.788 & Bid values & 0.0000 \\
Combined (FL + Imhof)$^{\ddagger}$ & 0.502 & Both & 0.0000 \\
Random classifier & 0.500 & None & --- \\
\bottomrule
\end{tabular}
\begin{tablenotes}
\small
\item \textit{Notes:} AUC computed on 16,692 always-loser firms
(98 CADE co-bidders, 16,594 non-CADE).
Ground truth: co-participation with CADE-convicted cartelists.
FL screen score: tender participation count (higher $=$ more suspicious).
Imhof composite: standardized mean of within-tender CV, excess kurtosis,
skewness, and normalized bid spread across all tenders the firm participates in.
DeLong $p$-value tests the null that the AUC difference from FL screen equals zero.
$^{\ddagger}$Combined score averages standardized FL and Imhof
$z$-scores. The near-random AUC is a direct consequence of the
dispersion paradox: FL firms (cover bidders) have \emph{lower}
bid dispersion than genuine losers, so the Imhof scores---which
flag \emph{high} dispersion---are negatively correlated with the
FL participation score. Averaging anti-correlated signals cancels
the detection power. This is not a coding artifact; it is the
empirical manifestation of Regime~2: any combined screen that
adds dispersion features to a participation-based screen will
\emph{degrade} performance when coordinated cover bidding
dominates. FL screen AUC differs slightly across tables
(0.937--0.940) due to different evaluation samples.
\end{tablenotes}
\end{threeparttable}
\end{table}
